% Shared preamble for the three ENS003 software-team presentations.
% Each presentation \input{}'s this file.

\documentclass[11pt,aspectratio=169]{beamer}

\usepackage[utf8]{inputenc}
\usepackage[T1]{fontenc}
\usepackage{lmodern}
\usepackage{amsmath}
\usepackage{booktabs}
\usepackage{graphicx}
\usepackage{listings}
\usepackage{xcolor}
\usepackage{hyperref}

% Built-in Beamer theme so the deck compiles on any TeX Live / MacTeX
% install without external packages. Customise colours below.
\usetheme{Boadilla}
\usecolortheme{seahorse}
\setbeamertemplate{navigation symbols}{}
\setbeamertemplate{footline}[frame number]

\definecolor{ink}{HTML}{1A1A1A}
\definecolor{accentred}{HTML}{A02828}
\definecolor{accentgreen}{HTML}{2C662C}
\definecolor{codebg}{HTML}{F7F5EE}
\definecolor{rulegrey}{HTML}{888888}

\setbeamercolor{normal text}{fg=ink}
\setbeamercolor{frametitle}{bg=ink, fg=white}
\setbeamercolor{title}{fg=ink}
\setbeamercolor{author}{fg=ink}
\setbeamercolor{date}{fg=rulegrey}
\setbeamercolor{block title}{bg=ink, fg=white}
\setbeamercolor{block body}{bg=codebg, fg=ink}
\setbeamercolor{alerted text}{fg=accentred}
\setbeamercolor{structure}{fg=ink}
\setbeamerfont{frametitle}{series=\bfseries}

\lstset{
  basicstyle=\ttfamily\scriptsize,
  backgroundcolor=\color{codebg},
  language=Python,
  keywordstyle=\color{accentred},
  stringstyle=\color{accentgreen},
  commentstyle=\color{rulegrey}\itshape,
  breaklines=true,
  numbers=none,
  showstringspaces=false,
  frame=single,
  framerule=0pt,
  framesep=6pt,
  xleftmargin=0pt,
  xrightmargin=0pt
}

% Convenience commands
\newcommand{\code}[1]{\texttt{\small #1}}
\newcommand{\file}[1]{\textcolor{accentred}{\texttt{\small #1}}}
\newcommand{\unit}[1]{\,\mathrm{#1}}

% Title-slide metadata is set per-deck (so the same preamble can power
% both Turkish and English versions of each talk).


\newcommand{\courseHeader}{ENS003 --- Digital Twins for Health Sciences}
\newcommand{\courseSubheader}{Istinye University · Faculty of Engineering and Natural Sciences}
\newcommand{\instructor}{Instructor: Prof. Dr. Şenol Pişkin}
\newcommand{\projectTitle}{A Surrogate ML Model for a Ti-6Al-4V Cementless Hip Implant}

\title{Data Pipeline and Feature Engineering}
\subtitle{\projectTitle}
\author{Ata Bulut \\ \small Computer Engineering · 220901582}
\date{May 2026 \\ \scriptsize \courseHeader{} · \instructor}

\begin{document}

\maketitle

% ---------------------------------------------------------------------------
\begin{frame}{In this talk}
  \begin{itemize}
    \item Turning the FEA team's raw output into something an ML model can eat
    \item The pipeline from ANSYS export to \file{data/dataset.csv}
    \item Architecture of \file{convert.py} and \file{preprocess.py}
    \item Bugs I hit, tests and coverage
  \end{itemize}
  \vspace{0.4cm}
  \begin{block}{My responsibility (software team)}
    Data conversion, schema design, feature engineering,
    and the CV grouping helper.
  \end{block}
\end{frame}

% ---------------------------------------------------------------------------
\begin{frame}{The data source}
  ANSYS Workbench parametric Design Point sweep:
  \begin{itemize}
    \item \(m \in \{50, 60, 70, 80, 90, 100\}\unit{kg}\) (6 values)
    \item \(K \in \{2, 2.5, \ldots, 6\}\) (9 values, body-weight multiplier)
    \item \(\theta \in \{0°, 10°, 20°\}\) (3 values, force tilt angle)
  \end{itemize}
  \vspace{0.3cm}
  \begin{center}
    \(\boldsymbol{6 \times 9 \times 3 = 162}\) unique design points
  \end{center}
  \vspace{0.3cm}
  ANSYS additionally re-ran 108 of them after a project update --- 270
  rows in total. Duplicates differ only at the 5th--6th significant figure
  (solver micro-noise).

  \vspace{0.3cm}
  \begin{block}{Call: keep all 270 rows}
    \footnotesize
    Every ANSYS solve was expensive --- throwing them away is wasteful.
    To avoid train/test leakage, cross-validation is
    \alert{grouped on \(\boldsymbol{(m, K, \theta)}\)} so duplicates
    always end up in the same fold.
  \end{block}
\end{frame}

% ---------------------------------------------------------------------------
\begin{frame}[fragile]{\file{convert.py} architecture}
\begin{lstlisting}
def convert(source=None, out=None):
    src = source if source else _pick_source(prefer_csv=True)
    if src.suffix.lower() == ".csv":
        raw = _read_csv(src)          # skip 7-line preamble
    else:
        raw = _read_xlsx(src)         # drop the "Units" row

    df = _normalize_columns(raw)      # "P1" or "P1 - Force..."
    df = pd.to_numeric(...).dropna()
    df = _convert_units(df)           # MPa -> Pa, mm -> m
    df = _attach_constants(df, ModelConstants())
    df["case_id"] = [f"DP_{i+1:03d}" for i in range(len(df))]
    df.to_csv(out, index=False)
    return df
\end{lstlisting}

  \begin{itemize}\small
    \item Two source formats: CSV (preamble + bare \code{P1}) and XLSX
          (Units row + \code{P1 - Force X Component})
    \item A single regex \verb|^\s*(P\d+)\b| catches both header forms
    \item Unit normalisation: MPa\(\to\)Pa, mm\(\to\)m
    \item ANSYS constants (mesh, material) attached row-wise from \code{ModelConstants}
  \end{itemize}
\end{frame}

% ---------------------------------------------------------------------------
\begin{frame}[fragile]{\file{preprocess.py}: feature engineering}
\begin{lstlisting}
INPUT_FEATURES = ["patient_mass_kg", "K_factor", "force_angle_deg"]
DERIVED_FEATURES = ["force_x_N", "force_z_N", "force_magnitude_N"]

def add_features(df):
    out = df.copy()
    rad = np.deg2rad(out["force_angle_deg"])
    Fmag = out["patient_mass_kg"] * 9.81 * out["K_factor"]
    out["force_x_N"] = Fmag * np.cos(rad)
    out["force_z_N"] = Fmag * np.sin(rad)
    out["force_magnitude_N"] = Fmag
    return out
\end{lstlisting}

  \vspace{0.2cm}
  \begin{itemize}\small
    \item The model sees three independent inputs: \(m, K, \theta\)
    \item \(F_x, F_z, |F|\) are derived \textbf{analytically} from those three
          --- predict-time and train-time use the same formula, no drift
    \item \(F_y\) and \code{fixed\_support\_n\_faces} are constant across all
          rows (zero variance) \(\Rightarrow\) dropped from the input vector
  \end{itemize}
\end{frame}

% ---------------------------------------------------------------------------
\begin{frame}[fragile]{CV grouping --- \code{grid\_groups()}}
\begin{lstlisting}
GRID_KEYS = ("patient_mass_kg", "K_factor", "force_angle_deg")

def grid_groups(df):
    """Return a group ID per row, equal for ANSYS re-runs."""
    grouped = df[list(GRID_KEYS)].astype(str).agg("|".join, axis=1)
    return grouped.to_numpy()
\end{lstlisting}

  \vspace{0.4cm}
  \textbf{How Recep wires it in \code{train.py}:}
  \begin{itemize}\small
    \item \code{GroupKFold(n\_splits=5).split(X, y, groups=grid\_groups(df))}
    \item Every duplicate of the same \((m, K, \theta)\) triplet shares a fold
    \item Without this, \(R^2\) would inflate to \(\approx 1.0\) (leakage)
  \end{itemize}
\end{frame}

% ---------------------------------------------------------------------------
\begin{frame}{Bugs I hit on the way}
  \begin{itemize}
    \item \textbf{The XLSX ``Units'' row:} the first data row is actually a
          unit label (\code{kg, MPa, mm, ...}). pandas inferred \code{object}
          dtype for the whole column, breaking arithmetic.
          \alert{Fix:} drop the row, then \code{pd.to\_numeric(errors=`coerce')}
    \item \textbf{Two header formats:} CSV has \code{P1}, XLSX has
          \code{P1 - Force X Component}. \alert{Fix:} a single regex
          \verb|^\s*(P\d+)\b| catches both.
    \item \textbf{\code{out=OUT\_CSV} default argument:} Python binds defaults
          at definition time, so \code{monkeypatch} in tests was a no-op ---
          tests silently wrote to the real \code{data/dataset.csv}.
          \alert{Fix:} \code{out=None} + resolve inside the function.
  \end{itemize}
\end{frame}

% ---------------------------------------------------------------------------
\begin{frame}{Tests \& coverage}
  \begin{center}
    \begin{tabular}{lrr}
      \toprule
      File & Statements & Coverage \\
      \midrule
      \file{src/convert.py}      & 115 & 90\% \\
      \file{src/preprocess.py}   & 27  & 100\% \\
      \bottomrule
    \end{tabular}
  \end{center}

  \vspace{0.5cm}
  \textbf{Key test cases:}
  \begin{itemize}\small
    \item Immutability: \code{add\_features()} does not mutate its input
    \item Trigonometry: \(F_x^2 + F_z^2 = |F|^2\) for every row
    \item Stale \(F_x\)/\(F_z\) columns get overwritten (predict-time consistency)
    \item ANSYS re-runs share a single \code{grid\_groups()} ID
  \end{itemize}
\end{frame}

% ---------------------------------------------------------------------------
\begin{frame}[fragile]{Demo}
\begin{lstlisting}[language=bash, basicstyle=\ttfamily\small]
$ ./run.sh convert
wrote 270 rows -> data/dataset.csv
total rows kept: 270 (unique grid points: 162)

$ ./run.sh check | head -6
rows=270  columns=30
unique grid points: 162
non-null target counts:
  safety_factor_neck_min: 270 / 270
  ...
\end{lstlisting}

  \vspace{0.4cm}
  The output feeds directly into Recep's \file{train.py}. The schema is
  the \alert{single source of truth} we share with the modelling side.
\end{frame}

% ---------------------------------------------------------------------------
\begin{frame}{Questions}
  \vfill
  \begin{center}
    \Large Thank you.\\[0.4cm]
    \normalsize Questions?
  \end{center}
  \vfill
  \scriptsize
  \begin{itemize}
    \item \file{src/convert.py}, \file{src/preprocess.py}
    \item \file{tests/test\_preprocess.py} (12 tests, 100\% coverage)
    \item Development log: \file{docs/PROJECT\_NOTES.md} Sprints 1--2
  \end{itemize}
\end{frame}

\end{document}
